\section{Ispezione del codice e analisi statica}
\label{sec:analisi-statica}

Questa parte l'ho affrontata in tre passaggi diversi, perché ognuno trova cose che gli
altri due si perdono: prima l'ispezione manuale con la checklist vista a lezione, poi
una seconda lettura fatta fare a un LLM, e infine l'analisi automatica con SpotBugs.

\subsection{Ispezione con la checklist}
\label{sec:checklist}

La prima passata l'ho fatta con la checklist di code inspection del corso, andando in
ordine sulle nove categorie. Sotto riporto per ognuna cosa è venuto fuori. Alcune voci
sul mio codice non si applicano proprio, e lo dico invece di spuntarle a vuoto.

\paragraph{1. Specification / Design.}
La funzionalità descritta nei requisiti c'è tutta: ogni requisito da RF1 a RF14 ha il
suo corrispettivo in \texttt{SmartParkingFSM} o in \texttt{Parcheggio}. Di funzionalità
in più rispetto alla specifica non ne ho trovate: i due metodi che sembravano candidati,
\texttt{SensorInput.nessunEvento()} e il costruttore \texttt{SmartParkingFSM(Display)},
sono entrambi usati dai test, il primo per provare i self-loop (uno step senza nessun
sensore attivo) e il secondo per costruire la FSM con il parcheggio di default.

\paragraph{2. Initialization and Declarations.}
Tutti i campi sono inizializzati nel costruttore, e quelli che non cambiano mai
(\texttt{parcheggio}, \texttt{display}, tutti i campi di \texttt{SensorInput}) sono
\texttt{final}. I tipi sono quelli giusti: gli stati e i tipi di utente sono enum e non
interi o stringhe, quindi un valore fuori dominio non è nemmeno rappresentabile. Tutti i
campi sono \texttt{private}, e \texttt{SensorInput} è dichiarata \texttt{final} e
immutabile.

\paragraph{3. Method Calls.}
Qui è uscita la segnalazione più concreta di tutta la checklist, e non l'aveva vista
nessun altro controllo. Ci sono due firme con \textbf{parametri adiacenti dello stesso
tipo}:

\begin{lstlisting}[language=Java, caption={Firme con parametri adiacenti dello stesso tipo}]
public Parcheggio(int postiStdIniziali, int postiDisIniziali)

public SensorInput(boolean sensIn, boolean sensOut, boolean transitoOk,
                   boolean autoVia, TipoUtente utenteRilevato, boolean pagamentoOk)
\end{lstlisting}

Se in una chiamata scambio i due \texttt{int}, o due dei quattro \texttt{boolean}, il
codice compila senza un fiato e il comportamento cambia. È esattamente il tipo di errore
che la categoria 3 della checklist va a cercare, e che né il compilatore né i contratti
JML possono intercettare, perché a livello di tipi la chiamata è corretta.

Su \texttt{SensorInput} il problema in pratica non si presenta, perché il costruttore non
lo chiama mai nessuno direttamente: si passa sempre dai metodi statici
(\texttt{sensIn()}, \texttt{transito()}, \texttt{pagamento(esito)}\dots), che hanno zero
o un parametro e un nome che dice cosa fanno. È una protezione che avevo messo per
comodità di scrittura dei test, e che a posteriori risolve anche questo.

Su \texttt{Parcheggio(int, int)} invece il rischio resta. L'ho lasciato così perché quel
costruttore lo usano solo i test, dove i valori sono scritti in chiaro riga per riga e un
eventuale scambio farebbe fallire subito l'assert. In un sistema vero la soluzione
sarebbe un tipo dedicato per i due contatori, o un costruttore con \emph{builder}.

\paragraph{4. Arrays.}
Non applicabile: nel progetto non c'è nessun array né nessuna collezione. Lo stato è
tenuto in due interi e in due enum. Di conseguenza non ci sono errori di indice
off-by-one né rischi di sforare i limiti.

\paragraph{5. Object Comparison.}
Tutti i confronti tra enum sono fatti con \texttt{==}, che per gli enum Java è corretto e
anzi preferibile a \texttt{equals}. Stringhe da confrontare non ce ne sono: le uniche che
compaiono sono i messaggi passati al display, che vengono solo scritti, mai testati per
uguaglianza nel codice di produzione.

\paragraph{6. Output Format.}
I messaggi del display sono coerenti e senza errori di battitura. Due osservazioni
minori. La prima è che alcune stringhe sono ripetute in più punti (``Sistema pronto''
compare tre volte, una per ogni ritorno in IDLE): andrebbero in costanti, ma su quattro
messaggi in croce ho lasciato stare. La seconda è che ``Sbarra aperta'' e ``Sbarra di
uscita aperta'' sembrano un'incoerenza e invece sono volute, perché distinguono le due
sbarre.

\paragraph{7. Computation, Comparisons and Assignments.}
Divisioni non ce ne sono, quindi il problema del denominatore a zero non si pone.
L'aritmetica è solo somma e sottrazione di 1 sui due contatori, e non può andare in
overflow perché i contatori sono limitati tra 0 e la capacità massima, cosa garantita in
tre modi diversi: dall'invariante JML, dai controlli a runtime in \texttt{liberaPosto} e
dalle CTLSPEC sul modello. Vale la pena dirlo perché l'overflow è il classico punto in
cui la verifica ESC si inceppa, e qui non succede proprio grazie a quei vincoli.

\paragraph{8. Exceptions.}
Nel progetto non c'è nemmeno un \texttt{try/catch}. È una scelta: le uniche eccezioni
lanciate (\texttt{IllegalArgumentException} e \texttt{IllegalStateException}) segnalano
errori di programmazione, non condizioni che possono capitare a runtime, quindi devono
propagare e non essere nascoste. Un punto di attenzione però c'è: nei gestori dei bottoni
della UI Vaadin non c'è nessuna gestione, quindi un'eccezione lì diventerebbe una pagina
di errore del framework. Per un prototipo va bene, in una versione vera ci vorrebbe un
messaggio all'utente.

\paragraph{9. Flow of Control.}
Lo \texttt{switch} di \texttt{step} ha tutti i \texttt{case} chiusi da \texttt{break} e
un ramo \texttt{default} che lancia eccezione, quindi niente fall-through accidentale.
Il problema del \emph{dangling else} non si presenta: l'unico punto con due condizioni in
fila, \texttt{gestisciIdle}, usa un \texttt{if/else if} esplicito. Infine, nel package
\texttt{core} non c'è \textbf{nessun ciclo}: la FSM avanza di un passo per chiamata e
l'iterazione, se serve, la fa il chiamante. Questo spiega anche perché nella parte JML
(Sezione~\ref{sec:jml-ref}) non compare nessuna \texttt{loop\_invariant}: non essendoci
loop, non c'è niente da invariantizzare.

\subsection{Seconda lettura con un LLM}
\label{sec:llm-review}

La consegna chiedeva di fare l'ispezione ``con chatgpt o LLM in generale'', quindi dopo
la checklist ho usato ChatGPT come se fosse un collega a cui far dare una seconda
occhiata al codice. Gli ho passato tutto il sorgente dei package \texttt{core} e
\texttt{ui} e gli ho chiesto di leggerlo cercando nomi poco chiari, casi limite non
gestiti, possibili \texttt{NullPointerException} e punti in cui il codice si allontana
dal modello ASM.

Le cose che sono uscite non le ho prese per buone così com'erano: le ho valutate una per
una, decidendo di volta in volta se correggerle, mitigarle o lasciarle stare spiegando
perché. L'LLM ha segnalato i punti su cui ragionare, le decisioni le ho prese io.

\begin{enumerate}
    \item \textbf{Manca il controllo di null su \texttt{utenteRilevato} in CHK\_OUT}: il
    confronto restituisce false senza dire niente, e il sistema va in USC come se
    l'utente non fosse STD. L'ho lasciata così per restare fedele all'ASM, dove il
    dominio \texttt{TipoUtente} non prevede l'assenza del dato; in più tutti i punti che
    chiamano quel metodo impostano sempre l'utente.

    \item \textbf{\texttt{TipoPosto.NESSUNO} vuol dire due cose diverse}: in
    \texttt{assegnaPosto} significa ``non c'è posto'', in \texttt{liberaPosto} significa
    ``non fare niente''. Chiarito nel Javadoc, invece di inventare un tipo separato che
    avrebbe appesantito una classe che doveva restare semplice.

    \item \textbf{\texttt{liberaPosto} lancia eccezione se il contatore è già al
    massimo}. Questa l'ho tenuta apposta: preferisco un crash che si vede subito a un
    contatore che va oltre la capacità senza dire niente, violando l'invariante. Il caso
    è comunque coperto da un test.

    \item \textbf{Il decremento avviene in un momento diverso rispetto all'ASM}: qui
    \texttt{assegnaPosto} decide e scala insieme, nell'ASM si scala al transito. Negli
    scenari di test la differenza non si vede, quindi resta la versione più semplice,
    scritta nel Javadoc della classe.

    \item \textbf{Non c'è thread-safety}. Va bene così, perché c'è un varco solo pilotato
    da una FSM sequenziale, esattamente come nel modello. Resta segnata come limite se un
    giorno si volessero gestire più varchi.

    \item \textbf{\texttt{posto\_assegnato} è una memoria da un posto solo, e il difetto
    c'è anche nel modello ASM}. Questa è l'osservazione più importante di tutte. Sia la
    FSM Java sia l'ASM si ricordano un solo posto assegnato, ma nel parcheggio ci stanno
    due veicoli, uno standard e uno disabile. Se entra uno STD, poi entra un DISABILE che
    sovrascrive lo slot, e poi esce il primo, viene rilasciato il posto sbagliato e i
    contatori si sballano.

    Il punto è che l'implementazione è fedele alla specifica: è la specifica stessa che
    dà per scontato di tracciare un veicolo alla volta, senza mai associare un veicolo al
    suo posto. Per sistemarlo servirebbe una mappa veicolo$\to$posto sia nell'ASM sia in
    Java. L'ho lasciato e documentato come limite noto, perché mi sembra l'esempio più
    utile di tutto il progetto: è un difetto che il testing di conformità non può
    trovare, visto che il codice fa esattamente quello che dice il modello, e che salta
    fuori solo leggendo il codice con occhio critico.

    \item \textbf{I bottoni della UI sono sempre attivi}: \texttt{MainView} non li
    disabilita in base allo stato, quindi si può cliccare ``Pagamento'' mentre siamo in
    IDLE. Non succede niente di male perché quegli step finiscono assorbiti dai
    self-loop della FSM, cioè proprio dagli \texttt{if} senza \texttt{else} che il Model
    Advisor segnalava come MP2 (\S\ref{sec:asmetav-vc}). Per un prototipo va bene; in una
    versione vera si metterebbe un \texttt{setEnabled} legato allo stato.
\end{enumerate}

\subsection{Bilancio}

Fra checklist e LLM sono uscite otto osservazioni. Nessun bug vero da correggere nel
codice, e in due casi mi è bastato chiarire le cose nel Javadoc.

Quello che mi ha colpito è che i due metodi hanno trovato cose diverse. La checklist,
proprio perché è meccanica e va per categorie, mi ha fatto guardare le \emph{firme} dei
metodi, che leggendo il codice per capire cosa fa non guardi mai, e da lì è saltato fuori
il problema dei parametri adiacenti dello stesso tipo. L'LLM, che invece ragiona sul
significato, ha trovato il difetto della memoria a slot singolo, che una checklist non
avrebbe mai potuto vedere perché non è una violazione di nessuna regola: è il modello a
essere incompleto.

E quel difetto, il punto 6, è la cosa più utile emersa da tutto il progetto, perché né il
testing di conformità né la verifica ESC riuscivano a vederlo. I contratti JML di
\texttt{Parcheggio} restano rispettati anche nell'esecuzione che rilascia il posto
sbagliato, dato che il problema sta più in alto, a livello di FSM e di modello. Il model
checking e l'ESC dimostrano quello che è stato formalizzato; l'ispezione a mano trova
quello che nella formalizzazione non c'è.
